% Threat model for the SaTML 2027 submission (feat-014 prep, drafted 2026-09-05; prose revised 2026-09-05). \input from satml_2027.tex.
\section{Threat Model and Audit Goals}\label{sec:threat}

\paragraph{The mechanism under audit}
Anchored Decoding~\cite{he2026anchored} composes a \emph{risky} language model $p_r$, which may have been trained on copyrighted works, with a \emph{safe} anchor $p_s$ trained only on permissively licensed and public-domain text~\cite{kandpal2025commonpile}; the anchor is TinyComma-1.8B, and the released logs audited here~\cite{vijayavallabh2026audit} pair it with Llama-3.1-8B-Instruct. At every decoding step the served distribution $p_\theta \propto p_s^{1-\theta} p_r^{\theta}$ is the point on the exponential-family geodesic between the two next-token distributions whose divergence from the anchor does not exceed a per-step allowance $k_t$. The allowances are banked: with a per-token rate $k$ and a prompt-dependent prefix debt $\delta_{\mathrm{init}}(x)$, step $t$ may spend $k_t = \max\{0, (t+1)k - \sum_{i<t} a_i - \delta_{\mathrm{init}}(x)\}$, where $a_i = D_{\mathrm{KL}}(p_\theta \,\|\, p_s)$ is the realised expenditure at step $i$. The mechanism's stated guarantee (He et al., Theorem~3.1) is sequence-level $K$-near-access-freeness in the KL sense of Vyas et al.~\cite{vyas2023provable}: for every prompt $x$ and every generation length $T \le T_{\max}$,
\begin{equation}
D_{\mathrm{KL}}\bigl(p_\theta(\cdot \mid x)\,\|\,p_s(\cdot \mid x)\bigr) \le K = k\,T_{\max}.
\label{eq:knaf}
\end{equation}

\paragraph{Parties}
The \emph{deployer} runs the decoder and publishes $K$ as a certificate of copyright risk. The \emph{rights-holder} cares about an event: that a served output reproduces a substantial portion of a specific work~\cite{cooper2025files,haviv2026separate}. The \emph{adversary} is a user of the deployed system who controls the prompt, may issue many queries, may adapt each to previous outputs, and may choose the temperature. In the strongest variant the adversary also chooses $p_r$, which Eq.~\eqref{eq:knaf} permits because it holds for \emph{any} risky model; we instantiate this with a model fine-tuned to memorise the targets (Section~\ref{sec:memorising}). The same measurements apply, with no adversarial choice, to a deployer whose risky model memorised the works in pre-training; Cooper et al.~\cite{cooper2025extracting} show this for Llama~3.1~70B, the risky model the mechanism's authors evaluate, and Section~\ref{sec:natural} audits that case with the 70B base model on the passages it memorised. The adversary never modifies the decoder, the anchor, $k$, or the accounting.

\paragraph{What the certificate promises}
Eq.~\eqref{eq:knaf} is a statement about one prompt and one expectation. By the chain rule for KL divergence its left-hand side equals $\mathbb{E}_{y \sim p_\theta}[\sum_t a_t] = \mathbb{E}[Z]$, and the decoder enforces $Z \le \max\{0, B\}$ on every trajectory by construction. The certificate therefore says nothing directly about (i)~a single output, whose realised log-likelihood ratio $L(y) = \log p_\theta(y \mid x) - \log p_s(y \mid x)$ may exceed $K$; (ii)~the probability of the rights-holder's event, which the data-processing inequality bounds only through $S(x) = -\log p_s(y^\star \mid x)$, the anchor's surprisal of the protected passage $y^\star$; or (iii)~a sequence of prompts, since near-access-freeness does not compose~\cite{cohen2025blameless}. The audit is organised around these three gaps.

\paragraph{Audit goals}
We separate the questions and give each a test that can fail. \emph{Accounting} (C1): does the implementation satisfy its own invariant, $Z \le \max\{0,B\}+\varepsilon$ per trajectory and $a_t \le k_t + \varepsilon$ per step? \emph{Regime} (C2): at which $k$ does the constraint act at all? \emph{Certificate strength} (C3): for the protected passages of a benchmark, how often is the implied cap on reproduction, $\max\{p : d(p \,\|\, e^{-S(x)}) \le K\}$, equal to one? \emph{Event-level and composed leakage} (C4, C5): how heavy are the tails of $L(y)$, and how much of a passage can an adaptive adversary reconstruct with every query within budget? Only C1 is a claim the mechanism makes; C2--C5 measure what the claim is worth. Five further questions concern what would make it worth more, or what the answers are relative to: whether they change when the risky model is the one the mechanism's authors evaluate and the memorisation comes from pre-training (C7); whether the same decoder can enforce the $\Delta_{\max}$ form, and what the KL form concentrates to when it cannot (C8, C9); whether a per-user budget and a cap on the bank bound what many compliant queries reveal (C10); whether the vacuity survives a different permissive anchor and a longer work (C11); and what each variant costs on ordinary prompts when something other than the risky model scores it (C12). Section~\ref{sec:protocol} gives each a failure mode, and Sections~\ref{sec:natural} and~\ref{sec:certificates} answer them.

\paragraph{Out of scope}
We do not assess legal questions, non-literal reproduction, or variants of the mechanism, and we do not attack the anchor's training data. The audit reads the decoder's logs, as a deployer or third-party auditor could; the adversary needs only query access.
